\documentclass{article}
\usepackage{amsmath, amsthm, amssymb}

ewtheorem{theorem}{Theorem}

ewtheorem{remark}{Remark}

ewtheorem{example}{Example}

\begin{document}

\qquad\qquad\qquad\qquad\qquad\qquad\qquad\qquad\qquad\qquad 6 \text{ Integration}
\begin{theorem} \text{(Comparison theorem) Suppose the functions } x \mapsto f(x) \text{ and } x \mapsto g(x)
\text{are defined on the interval } [a, \omega[ \text{ and integrable on any closed interval } [a, b] \subset
[a, \omega[.
\text{If}
\begin{equation*}
0 \leq f(x) \leq g(x)
\end{equation*}
\text{on } [a, \omega[, \text{ then convergence of the integral (\ref{eq:6.73}) implies convergence of (\ref{eq:6.72}) and}
\text{the inequality}
\begin{equation*}
\int_{a}^{\omega} f(x)dx \leq \int_{a}^{\omega} g(x)dx
\end{equation*}
\text{holds. Divergence of the integral (\ref{eq:6.72}) implies divergence of (\ref{eq:6.73}).}
\end{theorem}
\begin{proof}
\text{From the hypotheses of the theorem and the inequalities for proper Riemann}
\text{integrals we have}
\begin{equation*}
F(b) = \int_{a}^{b} f(x)dx \leq \int_{a}^{b} g(x)dx = G(b)
\end{equation*}
\text{for any } b \in [a, \omega[. \text{ Since both functions } F \text{ and } G \text{ are nondecreasing on } [a, \omega[, \text{ the}
\text{theorem follows from this inequality and Proposition 3.}
\end{proof}
\begin{remark} \text{If instead of satisfying the inequalities } 0 \leq f (x) \leq g(x) \text{ the functions } f
\text{and } g \text{ in the theorem are known to be nonnegative and of the same order as } x \to \omega,
x \in [a, \omega[, \text{ that is, there are positive constants } c_1 \text{ and } c_2 \text{ such that}
\begin{equation*}
c_1 f (x) \leq g(x) \leq c_2 f (x),
\end{equation*}
\text{then by the linearity of the improper integral and theorem just proved, in this case}
\text{we can conclude that the integrals (\ref{eq:6.72}) and (\ref{eq:6.73}) either both converge or both}
\text{diverge.}
\end{remark}
\begin{example} \text{The integral}
\begin{equation*}
\int_{1}^{+\infty} \frac{\sqrt{x}}{\sqrt{1+x^4}} dx
\end{equation*}
\text{converges, since}
\begin{equation*}
\frac{\sqrt{x}}{\sqrt{1+x^4}} \sim \frac{1}{x^{3/2}}
\end{equation*}
\text{as } x \to +\infty.
\end{example}
\begin{example} \text{The integral}
\begin{equation*}
\int_{1}^{+\infty} \frac{\cos x}{x^2} dx
\end{equation*}
\text{converges absolutely, because}
\[
\left|\frac{\cos x}{x^2}\right| \leq \frac{1}{x^2},
\]
\text{and the integral } \(\displaystyle \int_{1}^{+\infty} \frac{1}{x^2} dx\) \text{ converges.}
\end{example}

\end{document}